%This work is licensed under the
%Creative Commons Attribution-Share Alike 3.0 United States License.
%To view a copy of this license, visit
%http://creativecommons.org/licenses/by-sa/3.0/us/ or send a letter to
%Creative Commons,
%171 Second Street, Suite 300,
%San Francisco, California, 94105, USA.

\chapter{Ogg Vorbis}
\label{vorbis}
Ogg Vorbis is a set of Vorbis audio packets
in an \hyperref[ogg]{Ogg container}.
All of the fields within Ogg Vorbis are little-endian.

\begin{figure}[h]
\includegraphics{figures/vorbis/stream.pdf}
\end{figure}
\par
\noindent
The first Ogg page contains the identification packet.
The second Ogg packet contains the comment.
The third Ogg packet contains the setup.
Subsequent Ogg packets contain audio.

\clearpage

\subsection{the Identification Packet}
The first packet within a Vorbis stream is the \VAR{Identification} packet.
This contains the sample rate and number of channels.
Vorbis does not have a bits-per-sample field, as samples
are stored internally as floating point values and are
converted into a certain number of bits in the decoding process.
To find the total samples, use the \VAR{Granule Position} value
in the stream's final Ogg page.
\par
\noindent
{\relsize{-1}
\ALGORITHM{the first Ogg packet}{Vorbis identification information}
\SetKwData{CHANNELS}{channels}
\SetKwData{SAMPLERATE}{sample rate}
\SetKwData{BITRATEMAX}{bitrate maximum}
\SetKwData{BITRATENOM}{bitrate nominal}
\SetKwData{BITRATEMIN}{bitrate minimum}
\SetKwData{BLOCKSIZE}{block size}
\ASSERT $(\text{\READ 8 unsigned bits}) = 1$\tcc*[r]{packet type}
\ASSERT $(\text{\READ 6 bytes}) = \texttt{"vorbis"}$\;
\BlankLine
\ASSERT $(\text{\READ 32 unsigned bits}) = 0$\tcc*[r]{Vorbis version}
$\CHANNELS \leftarrow$ \READ 8 unsigned bits\;
\ASSERT $\CHANNELS > 0$\;
$\SAMPLERATE \leftarrow$ \READ 32 unsigned bits\;
\ASSERT $\SAMPLERATE > 0$\;
$\BITRATEMAX \leftarrow$ \READ 32 signed bits\;
$\BITRATENOM \leftarrow$ \READ 32 signed bits\;
$\BITRATEMIN \leftarrow$ \READ 32 signed bits\;
$\text{\BLOCKSIZE}_0 \leftarrow 2 ^ {(\text{\READ 4 unsigned bits)}}$\;
\ASSERT $\text{\BLOCKSIZE}_0$ \IN \texttt{[64, 128, 256, 512, 1024, 2048, 4096, 8192]}\;
$\text{\BLOCKSIZE}_1 \leftarrow 2 ^ {(\text{\READ 4 unsigned bits})}$\;
\ASSERT $\text{\BLOCKSIZE}_1$ \IN \texttt{[64, 128, 256, 512, 1024, 2048, 4096, 8192]}\;
\ASSERT $\text{\BLOCKSIZE}_0 \leq \text{\BLOCKSIZE}_1$\;
\ASSERT $\text{(\READ 1 unsigned bit)} = 1$\tcc*[r]{framing bit}
\BlankLine
\Return $\left\lbrace\begin{tabular}{l}
\CHANNELS \\
\SAMPLERATE \\
\BITRATEMAX \\
\BITRATENOM \\
\BITRATEMIN \\
\BLOCKSIZE \\
\end{tabular}\right.$\;
\EALGORITHM
}
\begin{figure}[h]
\includegraphics{figures/vorbis/identification.pdf}
\end{figure}

\clearpage

\subsubsection{Identification Packet Example}

\begin{figure}[h]
  \includegraphics{figures/vorbis/identification_example.pdf}
\end{figure}
{\relsize{-1}
\begin{tabular}{rl}
packet type : & \texttt{1} \\
vorbis : & \texttt{"vorbis"} \\
version : & \texttt{0} \\
channels : & \texttt{2} \\
sample rate : & \texttt{44100} Hz \\
maximum bitrate : & \texttt{0} \\
nominal bitrate : & \texttt{112000} \\
minimum bitrate : & \texttt{0} \\
$\text{block size}_0$ : & $2 ^ \texttt{8} = 256$ \\
$\text{block size}_1$ : & $2 ^ \texttt{11} = 2048$ \\
framing : & 1 \\
\end{tabular}
}

\clearpage

\subsection{the Comment Packet}
\label{vorbiscomment}
The second packet within a Vorbis stream is the Comment packet.
\par
\noindent
{\relsize{-1}
\ALGORITHM{the second Ogg packet}{Vorbis comment information}
\SetKwData{VENLENGTH}{vendor string length}
\SetKwData{VENSTRING}{vendor string}
\SetKwData{COMCOUNT}{commment string count}
\SetKwData{LEN}{comment string length}
\SetKwData{COMMENT}{comment string}
\ASSERT $(\text{\READ 8 unsigned bits}) = 3$\tcc*[r]{packet type}
\ASSERT $(\text{\READ 6 bytes}) = \texttt{"vorbis"}$\;
\BlankLine
$\VENLENGTH \leftarrow$ \READ 32 unsigned bits\;
$\VENSTRING \leftarrow$ \READ (\VENLENGTH) bytes as UTF-8 string\;
$\COMCOUNT \leftarrow$ \READ 32 unsigned bits\;
\For{$i \leftarrow 0$ \emph{\KwTo}\COMCOUNT}{
  $\text{\LEN}_i \leftarrow$ \READ 32 unsigned bits\;
  $\text{\COMMENT}_i \leftarrow$ \READ ($\text{\LEN}_i$) bytes as UTF-8 string\;
}
\ASSERT $\text{(\READ 1 unsigned bit)} = 1$\tcc*[r]{framing bit}
\BlankLine
\Return $\left\lbrace\begin{tabular}{l}
\VENSTRING \\
\COMMENT \\
\end{tabular}\right.$\;
\EALGORITHM
}

\begin{figure}[h]
\includegraphics{figures/vorbis/comment.pdf}
\end{figure}
\par
\noindent
Comment strings are \texttt{"KEY=value"} pairs
where \texttt{KEY} is an ASCII value in the range \texttt{0x20}
through \texttt{0x7D}, excluding \texttt{0x3D},
is case-insensitive and may occur in multiple comment strings.
\texttt{value} is a UTF-8 value.
\begin{table}[h]
{\relsize{-1}
  \begin{tabular}{rlrl}
    \texttt{ALBUM} & album name &
    \texttt{ARTIST} & artist name \\
    \texttt{CATALOG} & CD spine number &
    \texttt{COMPOSER} & the work's author \\
    \texttt{COMMENT} & a short comment &
    \texttt{CONDUCTOR} & performing ensemble's leader \\
    \texttt{COPYRIGHT} & copyright attribution &
    \texttt{DATE} & recording date \\
    \texttt{DISCNUMBER} & disc number for multi-volume work &
    \texttt{DISCTOTAL} & disc total for multi-volume work \\
    \texttt{GENRE} & a short music genre label &
    \texttt{ISRC} & ISRC number for the track \\
    \texttt{PERFORMER} & performer name, orchestra, actor, etc. &
    \texttt{PUBLISHER} & album's publisher \\
    \texttt{SOURCE MEDIUM} & CD, radio, cassette, vinyl LP, etc. &
    \texttt{TITLE} & track name \\
    \texttt{TRACKNUMBER} & track number &
    \texttt{TRACKTOTAL} & total number of tracks \\
  \end{tabular}
}
\end{table}

\clearpage

\subsubsection{Comment Packet Example}

\begin{figure}[h]
  \includegraphics{figures/vorbis/comment-example.pdf}
\end{figure}
\begin{table}[h]
\begin{tabular}{rl}
  vendor string : & \texttt{"Xiph.Org libVorbis I 20090709"} \\
  comment string count : & \texttt{2} \\
  $\text{comment string}_0$ : & \texttt{"title=Title"} \\
  $\text{comment string}_1$ : & \texttt{"artist=Artist"} \\
\end{tabular}
\end{table}

%% \subsection{the Setup Packet}

%% The third packet in the Vorbis stream is the Setup packet.
%% \par
%% \noindent
%% {\relsize{-1}
%% \ALGORITHM{the third Ogg packet}{Vorbis decoding information}
%% \SetKwData{CBCOUNT}{codebook count}
%% \SetKwData{CB}{codebook}
%% \SetKwData{VTCOUNT}{time count}
%% \SetKwData{FLCOUNT}{floor count}
%% \SetKwData{FLTYPE}{floor type}
%% \SetKwData{FLOOR}{floor}
%% \SetKwData{RESCOUNT}{residue count}
%% \SetKwData{RESTYPE}{residue type}
%% \SetKwData{RESIDUE}{residue}
%% \SetKwData{CHANNELS}{audio channels}
%% \SetKwData{MAPCOUNT}{mapping count}
%% \SetKwData{MAPCONF}{mapping configuration}
%% \SetKwData{MODECOUNT}{mode count}
%% \SetKwData{MODECONF}{mode configuration}
%% \SetKwFunction{ILOG}{ilog}
%% \ASSERT $(\text{\READ 8 unsigned bits}) = 5$\tcc*[r]{packet type}
%% \ASSERT $(\text{\READ 6 bytes}) = \texttt{"vorbis"}$\;
%% \BlankLine
%% $\CBCOUNT \leftarrow \text{(\READ 8 unsigned bits)} + 1$\tcc*[r]{codebooks}
%% \For{$i \leftarrow 0$ \emph{\KwTo}\CBCOUNT}{
%%   $\text{\CB}_i \leftarrow$ \hyperref[vorbis:codebooks]{read codebook}\;
%% }
%% \BlankLine
%% $\VTCOUNT \leftarrow \text{(\READ 6 unsigned bits)} + 1$\tcc*[r]{time domain transforms}
%% \For{$i \leftarrow 0$ \emph{\KwTo}\VTCOUNT}{
%%   \ASSERT $(\text{\READ 16 unsigned bits}) = 0$\;
%% }
%% \BlankLine
%% $\FLCOUNT \leftarrow \text{(\READ 6 unsigned bits) + 1}$\tcc*[r]{floors}
%% \For{$i \leftarrow 0$ \emph{\KwTo}\FLCOUNT}{
%%   $\text{\FLTYPE}_i \leftarrow$ \READ 16 unsigned bits\;
%%   \uIf{$0 \leq \text{\FLTYPE}_i \leq 1$}{
%%     $\text{\FLOOR}_i \leftarrow$ \hyperref[vorbis:floors]{read floor of type $\text{\FLTYPE}_i$}\;
%%   }
%%   \lElse{
%%     unsupported floor type error\;
%%   }
%% }
%% \BlankLine
%% $\RESCOUNT \leftarrow \text{(\READ 6 unsigned bits) + 1}$\tcc*[r]{residues}
%% \For{$i \leftarrow 0$ \emph{\KwTo}\RESCOUNT}{
%%   $\text{\RESTYPE}_i \leftarrow$ \READ 16 unsigned bits\;
%%   \uIf{$0 \leq \text{\RESTYPE}_i \leq 2$}{
%%     $\text{\RESIDUE}_i \leftarrow$ read residue of type $\text{\RESTYPE}_i$\;
%%   }
%%   \lElse{
%%     unsupported residue type error\;
%%   }
%% }
%% $\MAPCOUNT \leftarrow \text{(\READ 6 unsigned bits) + 1}$\tcc*[r]{mappings}
%% \For{$i \leftarrow 0$ \emph{\KwTo}\MAPCOUNT}{
%%   $\text{\MAPCONF}_i \leftarrow$ \hyperref[vorbis:mappings]{read mapping}\;
%% }
%% \BlankLine
%% $\MODECOUNT \leftarrow (\text{\READ 6 unsigned bits}) + 1$\tcc*[r]{modes}
%% \For{$i \leftarrow 0$ \emph{\KwTo}\MODECOUNT}{
%%   $\MODECONF \leftarrow$ \hyperref[vorbis:mode]{read mode}\;
%% }
%% \BlankLine
%% \ASSERT $(\text{\READ 1 unsigned bit}) = 1$\tcc*[r]{framing bit}
%% \BlankLine
%% \Return $\left\lbrace\begin{tabular}{l}
%% \CBCOUNT \\
%% \CB \\
%% \VTCOUNT \\
%% \FLCOUNT \\
%% \FLTYPE \\
%% \FLOOR \\
%% \RESCOUNT \\
%% \RESTYPE \\
%% \RESIDUE \\
%% \MAPCOUNT \\
%% \MAPCONF \\
%% \MODECOUNT \\
%% \MODECONF \\
%% \end{tabular}\right.$\;
%% \EALGORITHM
%% }

%% \clearpage

%% \begin{figure}[h]
%%   \includegraphics{figures/vorbis/setup_packet.pdf}
%% \end{figure}

%% \clearpage

%% \subsubsection{Codebooks}
%% \label{vorbis:codebooks}
%% {\relsize{-1}
%% \ALGORITHM{Vorbis setup packet data}{Vorbis codebook information}
%% \SetKwData{CBDIMENSIONS}{codebook dimensions}
%% \SetKwData{CBENTRIES}{codebook entries}
%% \SetKwData{ORDERED}{ordered}
%% \SetKwData{SPARSE}{sparse}
%% \SetKwData{LENGTH}{codeword length}
%% \SetKwData{CENTRY}{current entry}
%% \SetKwData{CLENGTH}{current length}
%% \SetKwData{NUMBER}{number}
%% \SetKwData{CBLOOKUPTYPE}{codebook lookup type}
%% \SetKwData{CBMINVALUE}{codecook minimum value}
%% \SetKwData{CBDELTAVALUE}{codebook delta value}
%% \SetKwData{CBVALUEBITS}{codebook value bits}
%% \SetKwData{CBSEQUENCEP}{codebook sequence p}
%% \SetKwData{CBLOOKUPVAL}{codebook lookup values}
%% \SetKwData{CBMULT}{codebook multiplicands}
%% \SetKwFunction{ILOG}{ilog}
%% \SetKwFunction{FLOAT}{float32\_unpack}
%% \SetKwFunction{LOOKUPVALS}{lookup1\_values}
%% \ASSERT $\text{(\READ 24 unsigned bits) = \texttt{0x564342}}$\tcc*[r]{sync pattern}
%% $\CBDIMENSIONS \leftarrow$ \READ 16 unsigned bits\;
%% $\CBENTRIES \leftarrow$ \READ 24 unsigned bits\;
%% $\ORDERED \leftarrow$ \READ 1 unsigned bit\;
%% \eIf{$\ORDERED = 0$}{
%%   $\SPARSE \leftarrow$ \READ 1 unsigned bit\;
%%   \For{$i \leftarrow 0$ \emph{\KwTo}\CBENTRIES}{
%%     \eIf{$\SPARSE = 1$}{
%%       \eIf{$(\textnormal{\READ 1 unsigned bit}) = 1$}{
%%         $\text{\LENGTH}_i \leftarrow (\text{\READ 5 unsigned bits}) + 1$\;
%%       }{
%%         $\text{\LENGTH}_i \leftarrow$ marked as unused\;
%%       }
%%     }{
%%       $\text{\LENGTH}_i \leftarrow (\text{\READ 5 unsigned bits}) + 1$\;
%%     }
%%   }
%% }{
%%   $\CENTRY \leftarrow 0$\;
%%   $\CLENGTH \leftarrow (\text{\READ 5 unsigned bits}) + 1$\;
%%   \Repeat{$\CENTRY = \CBENTRIES$}{
%%     $\NUMBER \leftarrow$ \READ $\ILOG(\CBENTRIES - \CENTRY)$ unsigned bits\;
%%     \ASSERT $(\CENTRY + \NUMBER) \leq \CBENTRIES$\;
%%     \For{$i \leftarrow 0$ \emph{\KwTo}\NUMBER}{
%%       $\text{\LENGTH}_{\CENTRY} \leftarrow \CLENGTH$\;
%%       $\CENTRY \leftarrow \CENTRY + 1$\;
%%     }
%%     $\CLENGTH \leftarrow \CLENGTH + 1$\;
%%   }
%% }
%% \BlankLine
%% $\CBLOOKUPTYPE \leftarrow$ \READ 4 unsigned bits\;
%% \uIf{$1 \leq \CBLOOKUPTYPE \leq 2$}{
%%   $\CBMINVALUE \leftarrow \FLOAT(\text{\READ 32 unsigned bits})$\;
%%   $\CBDELTAVALUE \leftarrow \FLOAT(\text{\READ 32 unsigned bits})$\;
%%   $\CBVALUEBITS \leftarrow (\text{\READ 4 unsigned bits}) + 1$\;
%%   $\CBSEQUENCEP \leftarrow$ \READ 1 unsigned bit\;
%%   \eIf{$\CBLOOKUPTYPE = 1$}{
%%     $\CBLOOKUPVAL \leftarrow \LOOKUPVALS(\CBENTRIES~,~\CBDIMENSIONS)$\;
%%   }{
%%     $\CBLOOKUPVAL \leftarrow \CBENTRIES \times \CBDIMENSIONS$\;
%%   }
%%   \For{$i \leftarrow 0$ \emph{\KwTo}\CBLOOKUPVAL}{
%%     $\text{\CBMULT}_i \leftarrow$ \READ $(\CBVALUEBITS)$ unsigned bits\;
%%   }
%% }
%% \ElseIf{$\CBLOOKUPTYPE > 2$}{
%%   invalid lookup type error\;
%% }
%% \EALGORITHM
%% }

%% \clearpage

%% {\relsize{-1}
%% \begin{algorithm}[H]
%%   \SetKwData{CBLOOKUPTYPE}{codebook lookup type}
%%   \SetKwData{LAST}{last}
%%   \SetKwData{IDIVISOR}{index divisor}
%%   \SetKwData{MULTOFFSET}{multiplicant offset}
%%   \SetKwData{CBDIMENSIONS}{codebook dimensions}
%%   \SetKwData{LOOKUPOFFSET}{lookup offset} %???
%%   \SetKwData{CBLOOKUPVAL}{codebook lookup values}
%%   \SetKwData{VECTOR}{vector element}
%%   \SetKwData{CBMULT}{codebook multiplicands}
%%   \SetKwData{CBDELTAVALUE}{codebook delta value}
%%   \SetKwData{CBMINVALUE}{codecook minimum value}
%%   \SetKwData{CBSEQUENCEP}{codebook sequence p}
%%   \DontPrintSemicolon
%%   \uIf{$\CBLOOKUPTYPE = 1$}{
%%     $\LAST = 0$\;
%%     $\IDIVISOR \leftarrow 1$\;
%%     \For{$i \leftarrow 0$ \emph{\KwTo}\CBDIMENSIONS}{
%%       $\MULTOFFSET \leftarrow \lfloor\LOOKUPOFFSET \div \IDIVISOR \rfloor \bmod \CBLOOKUPVAL$\;
%%       $\text{\VECTOR}_{i} \leftarrow \text{\CBMULT}_{\MULTOFFSET} \times \CBDELTAVALUE + \CBMINVALUE + \LAST$\;
%%       \If{$\CBSEQUENCEP = 1$}{
%%         $\LAST \leftarrow \text{\VECTOR}_i$\;
%%       }
%%       $\IDIVISOR \leftarrow \IDIVISOR \times \CBLOOKUPVAL$\;
%%     }
%%   }
%%   \ElseIf{$\CBLOOKUPTYPE = 2$}{
%%     $\LAST = 0$\;
%%     $\MULTOFFSET \leftarrow \LOOKUPOFFSET \times \CBDIMENSIONS$\;
%%     \For{$i \leftarrow 0$ \emph{\KwTo}\CBDIMENSIONS}{
%%       $\text{\VECTOR}_i \leftarrow \text{\CBMULT}_{\MULTOFFSET} \times \CBDELTAVALUE + \CBMINVALUE + \LAST$\;
%%       \If{$\CBSEQUENCEP = 1$}{
%%         $\LAST \leftarrow \text{\VECTOR}_i$\;
%%       }
%%       $\MULTOFFSET \leftarrow \MULTOFFSET + 1$\;
%%     }
%%   }
%%   \Return \VECTOR\;
%% \end{algorithm}
%% }

%% \clearpage

%% \subsubsection{Floors}
%% \label{vorbis:floors}

%% {\relsize{-1}
%% \ALGORITHM{Vorbis setup packet data, floor type, maximum codebook number}{Vorbis floor information}
%% \SetKwData{FLTYPE}{floor type}
%% \SetKwData{ZORDER}{floor 0 order}
%% \SetKwData{ZRATE}{floor 0 rate}
%% \SetKwData{ZBARKMAPSIZE}{floor 0 bark map size}
%% \SetKwData{ZAMPBITS}{floor 0 amplitude bits}
%% \SetKwData{ZAMPOFFSET}{floor 0 amplitude offset}
%% \SetKwData{ZNUMBOOKS}{floor 0 number of books}
%% \SetKwData{ZBOOKLIST}{floor 0 book list}
%% \SetKwData{MAXCODEBOOK}{maximum codebook number}
%% \SetKwData{OPARTITIONS}{floor 1 partitions}
%% \SetKwData{OPCLASSLIST}{floor 1 partition class list}
%% \SetKwData{OCLASSLIST}{floor 1 class subclasses}
%% \SetKwData{ODIMENSIONS}{floor 1 class dimensions}
%% \SetKwData{OMASTERBOOKS}{floor 1 class master books}
%% \SetKwData{OSUBCLASSBOOKS}{floor 1 subclass books}
%% \SetKwData{OMULT}{floor 1 multiplier}
%% \SetKwData{OXLIST}{floor 1 X list}
%% \SetKwData{OVALUES}{floor 1 values}
%% \SetKwData{MAXCLASS}{maximum class}
%% \SetKwData{RANGEBITS}{range bits}
%% \SetKwData{CURCLASSNUM}{current class number}
%% \eIf{$\FLTYPE = 0$}{
%%   \begin{tabular}{rcl}
%%     \ZORDER & $\leftarrow$ & \READ 8 unsigned bits \\
%%     \ZRATE & $\leftarrow$ & \READ 8 unsigned bits \\
%%     \ZBARKMAPSIZE & $\leftarrow$ & \READ 16 unsigned bits \\
%%     \ZAMPBITS & $\leftarrow$ & \READ 6 unsigned bits \\
%%     \ZAMPOFFSET & $\leftarrow$ & \READ 8 unsigned bits \\
%%     \ZNUMBOOKS & $\leftarrow$ & $\text{(\READ 4 unsigned bits)} + 1$ \\
%%   \end{tabular}\;
%%   \For{$i \leftarrow 0$ \emph{\KwTo}\ZNUMBOOKS}{
%%     $\text{\ZBOOKLIST}_i \leftarrow$ \READ 8 unsigned bits\;
%%     \ASSERT $\text{\ZBOOKLIST}_i \leq \MAXCODEBOOK$\;
%%   }
%% }{
%%   $\OPARTITIONS \leftarrow$ \READ 5 unsigned bits\;
%%   \For{$i \leftarrow 0$ \emph{\KwTo}\OPARTITIONS}{
%%     $\text{\OPCLASSLIST}_i \leftarrow$ \READ 4 unsigned bits\;
%%   }
%%   $\MAXCLASS \leftarrow$ maximum value in $\text{\OPCLASSLIST}$, or 0 if $\OPARTITIONS = 0$\;
%%   \For{$i \leftarrow 0$ \emph{\KwTo}\MAXCLASS}{
%%     $\text{\ODIMENSIONS}_i \leftarrow (\text{\READ 3 unsigned bits}) + 1$\;
%%     $\text{\OCLASSLIST}_i \leftarrow$ \READ 2 unsigned bits\;
%%     \If{$\text{\OCLASSLIST}_i \neq 0$}{
%%       $\text{\OMASTERBOOKS}_i \leftarrow$ \READ 8 unsigned bits\;
%%     }
%%     \For{$j \leftarrow 0$ \emph{\KwTo}$2 ^ {\text{\OCLASSLIST}_i} - 1$}{
%%       $\text{\OSUBCLASSBOOKS}_{i~j} \leftarrow (\text{\READ 8 unsigned bits} - 1)$\;
%%     }
%%   }
%%   $\OMULT \leftarrow (\text{\READ 2 unsigned bits}) + 1$\;
%%   $\RANGEBITS \leftarrow$ \READ 4 unsigned bits\;
%%   $\text{\OXLIST}_0 \leftarrow 0$\;
%%   $\text{\OXLIST}_1 \leftarrow 2 ^ {\RANGEBITS}$\;
%%   $\OVALUES \leftarrow 2$\;
%%   \For{$i \leftarrow 0$ \emph{\KwTo}\OPARTITIONS}{
%%     $\CURCLASSNUM \leftarrow \text{\OPCLASSLIST}_i$\;
%%     \For{$j \leftarrow 0$ \emph{\KwTo}$\text{\ODIMENSIONS}_{\CURCLASSNUM}$}{
%%       $\text{\OXLIST}_{\OVALUES} \leftarrow$ \READ $(\RANGEBITS)$ unsigned bits\;
%%       $\OVALUES \leftarrow \OVALUES + 1$\;
%%     }
%%   }
%% }
%% \EALGORITHM
%% }

%% \clearpage

%% \subsubsection{Mappings}
%% \label{vorbis:mappings}
%% {\relsize{-1}
%% \ALGORITHM{Vorbis setup packet data}{Vorbis mapping information}
%% \SetKwData{CHANNELS}{audio channels}
%% \SetKwData{SUBMAPS}{mapping submaps}
%% \SetKwData{MAPCOUPSTEPS}{mapping coupling steps}
%% \SetKwData{MAPMAG}{mapping magnitude}
%% \SetKwData{MAPANGLE}{mapping angle}
%% \SetKwData{MAPMUX}{mapping mux}
%% \SetKwData{MAPSUBFLOOR}{mapping submap floor}
%% \SetKwData{MAPSUBRES}{mapping submap residue}
%% \ASSERT $(\text{\READ 16 unsigned bits}) = 0$\tcc*[r]{mapping type}
%% \eIf{$(\textnormal{\READ 1 unsigned bit}) = 1$}{
%%   $\text{\SUBMAPS} \leftarrow \text{(\READ 4 unsigned bits) + 1}$\;
%% }{
%%   $\text{\SUBMAPS} \leftarrow 1$\;
%% }
%% \eIf{$(\textnormal{\READ 1 unsigned bit}) = 1$}{
%%   $\text{\MAPCOUPSTEPS} \leftarrow (\text{\READ 8 unsigned bits}) + 1$\;
%%   \For{$j \leftarrow 0$ \emph{\KwTo}$\text{\MAPCOUPSTEPS}_i$}{
%%     $\text{\MAPMAG}_{j} \leftarrow$ \READ $\ILOG(\CHANNELS - 1)$ unsigned bits\;
%%     $\text{\MAPANGLE}_{j} \leftarrow$ \READ $\ILOG(\CHANNELS - 1)$ unsigned bits\;
%%     %%FIXME - add assertion here
%%   }
%% }{
%%   $\text{\MAPCOUPSTEPS} \leftarrow 0$\;
%% }
%% \ASSERT $(\text{\READ 2 unsigned bits}) = 0$\tcc*[r]{reserved}
%% \If{$\text{\SUBMAPS}  > 1$}{
%%   \For{$j \leftarrow 0$ \emph{\KwTo}\CHANNELS}{
%%     $\text{\MAPMUX}_{j} \leftarrow$ \READ 4 unsigned bits\;
%%     \ASSERT $\text{\MAPMUX}_{j} < \text{\SUBMAPS}_i$\;
%%   }
%% }
%% \For{$j \leftarrow 0$ \emph{\KwTo}$\text{\SUBMAPS}$}{
%%   \SKIP 8 bits\tcc*[r]{unused time configuration placeholder}
%%   $\text{\MAPSUBFLOOR}_{j} \leftarrow$ \READ 8 unsigned bits\;
%%   \ASSERT $\text{\MAPSUBFLOOR}_{j} \leq \text{highest number floor configurated for bitstream}$\;
%%   $\text{\MAPSUBRES}_{j} \leftarrow$ \READ 8 unsigned bits\;
%%   \ASSERT $\text{\MAPSUBRES}_{j} \leq \text{highest number residue configured for bitstream}$\;
%% }
%% \Return $\left\lbrace\begin{tabular}{l}
%% \SUBMAPS \\
%% \MAPCOUPSTEPS \\
%% \MAPMAG \\
%% \MAPANGLE \\
%% \MAPMUX \\
%% \MAPSUBFLOOR \\
%% \MAPSUBRES \\
%% \end{tabular}\right.$\;
%% \EALGORITHM
%% }

%% \clearpage

%% \subsubsection{Modes}
%% \label{vorbis:mode}
%% {\relsize{-1}
%% \ALGORITHM{Vorbis setup packet data}{Vorbis modes information}
%% \SetKwData{MODEBLOCKFLAG}{mode block flag}
%% \SetKwData{MODEWINDOWTYPE}{mode window type}
%% \SetKwData{MODETRANSFORMTYPE}{mode transform type}
%% \SetKwData{MODEMAPPING}{mode mapping}
%% $\text{\MODEBLOCKFLAG} \leftarrow$ \READ 1 unsigned bit\;
%% $\text{\MODEWINDOWTYPE} \leftarrow$ \READ 16 unsigned bits\;
%% \ASSERT $\text{\MODEWINDOWTYPE} = 0$\;
%% $\text{\MODETRANSFORMTYPE} \leftarrow$ \READ 16 unsigned bits\;
%% \ASSERT $\text{\MODETRANSFORMTYPE} = 0$\;
%% $\text{\MODEMAPPING} \leftarrow$ \READ 8 unsigned bits\;
%% \ASSERT $\text{\MODEMAPPING} \leq \text{highest number mapping in use}$\;
%% \Return $\left\lbrace\begin{tabular}{l}
%% \MODEBLOCKFLAG \\
%% \MODEWINDOWTYPE \\
%% \MODETRANSFORMTYPE \\
%% \MODEMAPPING \\
%% \end{tabular}\right.$\;
%% \EALGORITHM
%% }

%% \clearpage

%% \subsubsection{Codebooks}

%% The \VAR{Codebooks} section of the setup packet stores
%% the Huffman lookup trees.

%% \begin{figure}[h]
%% \includegraphics{figures/vorbis/codebooks.pdf}
%% \end{figure}
%% \par
%% \noindent
%% This section contains two optional sets of data,
%% a list of Huffman table entry lengths
%% and the lookup table values each entry length may resolve to.
%% \VAR{Total Entries} indicates the total number of entry lengths present.
%% These lengths may be stored in one of three ways:
%% unordered without sparse entries, unordered with sparse entries
%% and ordered.

%% Unordered without sparse entries is the simplest method;
%% each entry length is stored as a 5 bit value, plus one.
%% Unordered with sparse entries is almost as simple;
%% each 5 bit length is prefixed by a single bit indicating
%% whether it is present or not.

%% Ordered entries are more complicated.
%% The initial length is stored as a 5 bit value, plus one.
%% The number of entries with that length are stored as a series of
%% \VAR{Length Count} values in the bitstream, whose sizes
%% are determined by the number of remaining entries.
%% \begin{align*}
%% \text{Length Count}_i \text{ Size} &= \lfloor\log_2 (\text{Remaining Entries}_i)\rfloor + 1
%% \intertext{For example, given a \VAR{Total Entries} value of 8 and an
%% \VAR{Initial Length} value of 2:}
%% \text{Length Count}_0 \text{ Size} &= \lfloor\log_2 8\rfloor + 1 = 4 \text{ bits}
%% \end{align*}
%% which means we read a 4 bit value to determine how many
%% \VAR{Entry Length} values are 2 bits long.
%% Therefore, if $\text{Length Count}_0$ is read from the bitstream as 2,
%% our \VAR{Entry Length} values are:
%% \begin{align*}
%% \text{Entry Length}_0 &= 2 \\
%% \text{Entry Length}_1 &= 2
%% \end{align*}
%% and the next set of lengths are 3 bits long.
%% Since we still have remaining entries, we read another \VAR{Length Count}
%% value of the length:
%% \begin{equation*}
%% \text{Length Count}_1 \text{ Size} = \lfloor\log_2 6\rfloor + 1 = 3 \text{ bits}
%% \end{equation*}
%% Thus, if $\text{Length Count}_1$ is also a value of 2 from the bitstream,
%% our \VAR{Entry Length} values are:
%% \begin{align*}
%% \text{Entry Length}_2 &= 3 \\
%% \text{Entry Length}_3 &= 3
%% \end{align*}
%% and the next set of lengths are 4 bits long.
%% We then read one more \VAR{Length Count} value:
%% \begin{equation*}
%% \text{Length Count}_2 \text{ Size} = \lfloor\log_2 4\rfloor + 1 = 3 \text{ bits}
%% \end{equation*}
%% Finally, if $\text{Length Count}_2$ is 4 from the bitstream,
%% our \VAR{Entry Length} values are:
%% \begin{align*}
%% \text{Entry Length}_4 &= 4 \\
%% \text{Entry Length}_5 &= 4 \\
%% \text{Entry Length}_6 &= 4 \\
%% \text{Entry Length}_7 &= 4
%% \end{align*}
%% At this point, we've assigned lengths to all the values
%% indicated by \VAR{Total Entries}, so the process is complete.

%% \clearpage

%% \subsubsection{Transforming Entry Lengths to Huffman Tree}

%% Once a set of entry length values is parsed from the stream,
%% transforming them into a Huffman decision tree
%% is performed by taking our entry lengths from
%% $\text{Entry Length}_0$ to $\text{Entry Length}_{total - 1}$
%% and placing them in the tree recursively such that the 0 bit
%% branches are populated first.
%% For example, given the parsed entry length values:
%% \begin{align*}
%% \text{Entry Length}_0 &= 2 & \text{Entry Length}_1 &= 4 & \text{Entry Length}_2 &= 4 & \text{Entry Length}_3 &= 4 \\
%% \text{Entry Length}_4 &= 4 & \text{Entry Length}_5 &= 2 & \text{Entry Length}_6 &= 3 & \text{Entry Length}_7 &= 3
%% \end{align*}
%% \par
%% \noindent
%% We first place $\text{Entry Length}_0$ into the Huffman tree as a 2 bit value.
%% Since the zero bits are filled first when adding a node 2 bits deep,
%% it initially looks like:

%% \begin{figure}[h]
%% \includegraphics{figures/vorbis/huffman_example1.pdf}
%% \caption{$\text{Entry Length}_0$ placed with 2 bits}
%% \end{figure}
%% \par
%% \noindent
%% We then place $\text{Entry Length}_1$ into the Huffman tree as a 4 bit value.
%% Since the \texttt{0 0} branch is already populated by $\text{Entry Length}_0$,
%% we switch to the empty \texttt{0 1} branch as follows:
%% \begin{figure}[h]
%% \includegraphics{figures/vorbis/huffman_example2.pdf}
%% \caption{$\text{Entry Length}_1$ placed with 4 bits}
%% \end{figure}
%% \par
%% \noindent
%% The 4 bit $\text{Entry Length}_2$, $\text{Entry Length}_3$ and $\text{Entry Length}_4$
%% values are placed similarly along the \texttt{0 1} branch:
%% \begin{figure}[h]
%% \includegraphics{figures/vorbis/huffman_example3.pdf}
%% \caption{$\text{Entry Length}_2$, $\text{Entry Length}_3$ and $\text{Entry Length}_4$ placed with 4 bits}
%% \end{figure}
%% \par
%% \noindent
%% Finally, the remaining three entries are populated along the \texttt{1} branch:
%% \begin{figure}[h]
%% \includegraphics{figures/vorbis/huffman_example4.pdf}
%% \caption{$\text{Entry Length}_5$, $\text{Entry Length}_6$ and $\text{Entry Length}_7$ are placed}
%% \end{figure}

%% \subsubsection{The Lookup Table}

%% The lookup table is only present if \VAR{Lookup Type} is 1 or 2.
%% A \VAR{Lookup Type} of 0 indicates no lookup table, while anything
%% greater than 2 is an error.

%% \VAR{Minimum Value} and \VAR{Delta Value} are 32-bit floating point values
%% which can be parsed in the following way:
%% \begin{figure}[h]
%% \includegraphics{figures/vorbis/float32.pdf}
%% \end{figure}
%% \begin{equation*}
%% \text{Float} =
%% \begin{cases}
%% \text{Mantissa} \times 2 ^ {\text{Exponent} - 788} & \text{ if Sign = 0} \\
%% -\text{Mantissa} \times 2 ^ {\text{Exponent} - 788} & \text{ if Sign = 1}
%% \end{cases}
%% \end{equation*}
%% \par
%% \VAR{Value Bits} indicates the size of each value in bits, plus one.
%% \VAR{Sequence P} is a 1 bit flag.
%% The total number of multiplicand values depends on whether \VAR{Lookup Type}
%% is 1 or 2.
%% \begin{align*}
%% \intertext{if \VAR{Lookup Type} = 1:}
%% \text{Multiplicand Count} &= \text{max}(\text{Int}) ^ \text{Dimensions} \text{ where } \text{Int} ^ \text{Dimensions} \leq \text{Total Entries}
%% \intertext{if \VAR{Lookup Type} = 2:}
%% \text{Multiplicand Count} &= \text{Dimensions} \times \text{Total Entries}
%% \end{align*}
%% \par
%% The \VAR{Multiplicand} values themselves are a list of unsigned integers,
%% each \VAR{Value Bits} (+ 1) bits large.


\section{Channel Assignment}
\begin{table}[h]
{\relsize{-1}
\begin{tabular}{|c|r|r|r|r|r|r|r|r|}
\hline
channel & & & & & & & & \\
count & channel 1 & channel 2 & channel 3 & channel 4 & channel 5 & channel 6 & channel 7 & channel 8 \\
\hline
\multirow{2}{1em}{1} & front & & & & & & & \\
                     & center & & & & & & & \\
\hline
\multirow{2}{1em}{2} & front & front & & & & & & \\
                     & left  & right & & & & & & \\
\hline
\multirow{2}{1em}{3} & front & front & front & & & & & \\
                     & left  & center & right & & & & & \\
\hline
\multirow{2}{1em}{4} & front & front & back & back & & & & \\
                     & left  & right & left & right & & & & \\
\hline
\multirow{2}{1em}{5} & front & front  & front & back & back & & & \\
                     & left  & center & right & left & right & & & \\
\hline
\multirow{2}{1em}{6} & front & front  & front & back & back & & & \\
                     & left  & center & right & left & right & LFE & & \\
\hline
\multirow{2}{1em}{7} & front & front  & front & side & side  & back   & & \\
                     & left  & center & right & left & right & center & LFE& \\
\hline
\multirow{2}{1em}{8} & front & front  & front & side & side  & back   & back & \\
                     & left  & center & right & left & right & left & right & LFE \\
\hline
8+ & \multicolumn{8}{c|}{defined by application} \\
\hline
\end{tabular}
}
\end{table}
